\documentclass[11pt,a4paper]{article}
\usepackage[margin=1in]{geometry}
\usepackage{hyperref}
\usepackage{enumitem}
\usepackage{graphicx}
\usepackage{xcolor}
\usepackage{booktabs}
\usepackage{array}

% -----------------------------------------------------
% pkaushal-tiet-ucs503p-article.sty
% -----------------------------------------------------

% Variable: Lab Instructor
\makeatletter \newcommand{\BvrSetLabInstructor}[1]{%
\gdef\Bvr@LabInstructor{#1}}
% default
\newcommand{\Bvr@LabInstructor}{Nisha Thakur}
% public accessor
\newcommand{\BvrLabInstructorValue}{\Bvr@LabInstructor}
\makeatother

% Add subtitle
\usepackage{titling}
% -- Set title bold
\pretitle{\begin{center}\bfseries\fontsize{18bp}{18bp}\selectfont}
\makeatletter
\newcommand{\BvrSubtitle}[1]{%
\posttitle{%
\par\end{center}
\begin{center}\large#1\end{center}
\vskip 0.5em}%
}
\makeatother

% Hints in the section title(s)
\newcommand{\BvrHint}[1]{%
{\normalfont\normalsize\itshape\hfill #1}}

% -----------------------------------------------------

\title{Silo --- Collaborative Platform for Small Software}

\BvrSetLabInstructor{Nisha Thakur}

\BvrSubtitle{%
Project Proposal for UCS503P \\
Submitted to: \BvrLabInstructorValue \\ \bfseries
Thapar Institute of Engineering and Technology%
}
\author{
Paras Kaushal, CSED%
\thanks{Roll No: \texttt{1024030984},
Email: \texttt{pkaushal\_be24@thapar.edu}}
\and
Dev Kapla, CSED%
\thanks{Roll No: \texttt{1024030888},
Email: \texttt{dkapla\_be24@thapar.edu}}
\and
Devanshu, CSED%
\thanks{Roll No: \texttt{1024030876},
Email: \texttt{ddevanshu\_be24@thapar.edu}}
}

\date{August 18, 2026}

\begin{document}
\maketitle


% -------------------------------------------------------
\section{Higher-order goal%
\BvrHint{\ldots making small software a collaborative artifact}}

AI-assisted development has made it possible to create useful software for tasks that
previously would not justify a dedicated engineering project. A person can now describe a
dashboard, calculator, reporting tool, or internal workflow utility and obtain functioning
application code in minutes.

The problem is that producing the code is only the beginning. The resulting application
still needs to be configured, built, deployed, hosted, shared, updated, and maintained.
These steps typically require knowledge of programming environments, dependency management,
servers, deployment platforms, authentication, and version control.

This project proposes \textbf{Silo}, a collaborative platform that treats small software
more like a document than a traditional software project. A user should be able to create
or upload a small application without preparing deployment configuration; publish it as a
live application without understanding its underlying code; share it with others using
simple permissions; let others use, remix, or improve it; see what changed between versions
without reading source-code diffs; propose improvements without immediately changing the
version that others are using; and roll back when necessary.

The higher-order goal is therefore not simply to automate deployment. It is to make
\textbf{small, bespoke software a collaborative artifact}---something that can be created,
shared, improved, and governed with the same ease that people currently associate with
documents.

% -------------------------------------------------------
\section{Time-to-value %
\BvrHint{\ldots primary heuristic}}
\begin{itemize}[leftmargin=0.7cm]
  \item We will prioritise the shortest path from application code to a useful shared
    application: \textbf{Upload $\to$ Build $\to$ Publish $\to$ Share $\to$ Use
    $\to$ Remix $\to$ Review $\to$ Promote.}
  \item We will deliver meaningful increments quickly. The first iteration establishes
    the core loop; subsequent iterations add the collaborative governance layer on top
    of a working deployment.
  \item We will prioritise fast feedback over perfection by using weekly iterations
    with CI-backed automation from week~1 to reduce release friction.
  \item Success is defined by what can be delivered and measured in the first iteration,
    then improved based on user feedback. Specifically: a non-technical user uploads
    AI-generated code, receives a live URL, and shares it with a colleague---without
    either person writing a line of configuration.
\end{itemize}

% -------------------------------------------------------
\section{Problem Statement}

AI-assisted coding has reduced the cost of producing software. A person who needs a small
application can often generate one without formal software-development expertise. However,
the surrounding workflow remains designed for conventional software projects. A generated
application may contain several files but no deployment configuration, no reliable
dependency specification, no declared runtime, no known entry point, and no versioning
workflow suitable for non-developers. Consequently, many useful applications remain local
artefacts on a developer's laptop. The problem is not only deployment. There is also a
\emph{collaboration gap}:

\begin{itemize}[leftmargin=0.7cm]
  \item \textbf{No deployment path:} A user who can prompt an AI cannot write a
    Dockerfile or configure a reverse proxy. The app dies on the laptop it was born on.
  \item \textbf{No sharing model:} Sharing an application means sending source files,
    configuring a repository, and explaining how to run it---none of which are realistic
    for the people who benefit most from the tool.
  \item \textbf{No comprehensible change review:} When a colleague proposes an
    improvement, the owner is shown a source-code diff.
    A change such as \texttt{+~requests.post(\ldots)} does not communicate its practical
    significance. A more useful review would say:
    \emph{Added: sends data to \texttt{api.example.com}.
    Change type: capability-expanding.}
  \item \textbf{No safe update model:} There is no structured way to submit an
    improvement, review what changed in terms of what the application can now do, and
    publish the new version without unexpectedly breaking the version colleagues are
    already using.
  \item \textbf{No governance:} There is no audit trail, no revocation mechanism, and
    no administrator view of what internal tools exist or what they can access.
\end{itemize}

\textbf{Why this matters now:} AI-generated internal tooling is already proliferating
informally inside organisations. The platform to make it collaborative and governable does
not exist. Silo sits between AI-assisted software creation and everyday software use.

% -------------------------------------------------------
\section{Proposed Solution %
\BvrHint{\ldots Software Proposition}}

\subsection{Overview}

Build \textbf{Silo}---a web platform for publishing, sharing, and collaboratively evolving
small applications. The central product metaphor is \textbf{Google Docs for small
software.} The system combines four major capabilities:

\begin{itemize}[leftmargin=0.7cm]
  \item \textbf{Application publishing:} convert an application bundle into a working,
    hosted application with zero manual configuration.
  \item \textbf{Application understanding:} analyse the application and present a
    human-readable capability summary of what it requires and what it can access.
  \item \textbf{Software collaboration:} share applications using simple roles
    (Viewer\,/\,User\,/\,Forker\,/\,Editor\, /\,Owner); share via link exactly as one
    would share a document.
  \item \textbf{Capability-aware versioning:} allow applications to be remixed and
    updated while showing owners the meaningful capability changes between versions
    before a new version becomes the shared one.
\end{itemize}

The product has three surfaces: (1)~\textbf{Owner view}---upload code, receive a live
URL and a plain-English capability label, share via a dialog identical in feel to sharing
a Google Doc; (2)~\textbf{Recipient view}---open the link, see the capability summary,
use the application, optionally remix or propose an edit;
(3)~\textbf{Admin/org view}---all live applications, current permissions, pending reviews,
and one-click revoke.

\subsection{Core workflow}
\begin{enumerate}[leftmargin=0.7cm]
  \item User uploads AI-generated code (file or zip). Silo infers the runtime, resolves
    dependencies, locates the entry point, and deploys to a live subdomain URL
    automatically.
  \item Static analysis runs. A capability manifest is produced and a plain-English
    nutrition label is generated from it.
  \item Owner reviews the label, approves, and shares a link with a role.
  \item Recipient opens the link, sees the capability summary, and uses the running
    application.
  \item A forker or editor uploads an improved version. Silo produces the
    \emph{capability diff}---what this version can do that the current approved version
    cannot. The owner approves or rejects the diff, not the code.
  \item If approved, the stable application URL is atomically promoted to the new
    version. The previous version remains available for rollback.
\end{enumerate}

\subsection{Operational constraints for an educational setting}
\begin{itemize}[leftmargin=0.7cm]
  \item Single-VM deployment (Hetzner\,/\,Oracle free tier)---cost-free throughout.
  \item No proprietary cloud dependencies; all components self-hostable.
  \item Evaluation against a hand-built set of 10 representative applications generated
    using AI coding tools, covering diverse use cases.
\end{itemize}

% -------------------------------------------------------
\section{Solution Approach %
\BvrHint{\ldots Engineering focus}}

This is an engineering course project. The contribution is a novel composition of mature
components into a deployment-and-collaboration layer that does not currently exist,
evaluated against a representative set of AI-generated applications.

\subsection{Build inference engine %
\BvrHint{\ldots non-research}}

The build inference engine accepts an application bundle with no deployment configuration
and determines everything required to run it:

\begin{itemize}[leftmargin=0.7cm]
  \item Detect language and runtime version from syntax
    (e.g.\ \texttt{match}/\texttt{case} $\Rightarrow$ Python~$\geq$~3.10).
  \item Resolve import names to installable package names
    (\texttt{cv2}~$\to$~\texttt{opencv-python},
     \texttt{PIL}~$\to$~\texttt{Pillow},
     \texttt{sklearn}~$\to$~\texttt{scikit-learn}).
  \item Locate the entry point and bound port, including auto-repair of the common
    \texttt{127.0.0.1} binding that makes AI-generated apps unreachable inside a
    container.
  \item On build failure, apply bounded deterministic repair (max 3~attempts) and
    surface structured diagnostics.
\end{itemize}

The same tree-sitter parse that identifies the entry point provides the import list for
the capability analyser, so the two components share a single parse pass.

\subsection{Capability analysis and nutrition label %
\BvrHint{\ldots non-research}}

Capability detection is entirely deterministic---the LLM never decides safety:

\begin{itemize}[leftmargin=0.7cm]
  \item Semgrep rules target known sink patterns: socket calls, \texttt{open()},
    \texttt{subprocess}, \texttt{os.environ}, \texttt{requests}, \texttt{fetch}.
  \item Findings are normalised to canonical capability tuples. The \texttt{sha256}
    of the sorted tuple set is the manifest digest---a signed fingerprint of what this
    version can do.
  \item An LLM call converts the manifest to a plain-English nutrition label. Its only
    role is explanation; it cannot alter the findings.
\end{itemize}

\subsection{Capability diff and version approval %
\BvrHint{\ldots non-research}}

When a new version is submitted, its manifest is diffed against the currently approved
manifest. Each changed capability is classified as \textbf{neutral}, \textbf{reducing}
(fewer permissions), or \textbf{expanding} (new permissions). The owner reviews only
the diff---what this version is asking for that the current version does not have---and
approves or rejects it. The unit of review is the permission change, not the code.

\subsection{Web architecture %
\BvrHint{\ldots web-based solution}}

A standard 3-tier architecture:

\begin{itemize}[leftmargin=0.7cm]
  \item \textbf{Frontend:} Next.js + TypeScript + Tailwind---application workspace,
    nutrition-label/approval view, share dialog, version history, review interface,
    org dashboard.
  \item \textbf{Backend API:} FastAPI (Python) + Pydantic---build orchestration,
    analysis pipeline, version management, ACL enforcement, audit logging.
  \item \textbf{Database:} PostgreSQL---applications, versions, manifests,
    capabilities, approvals, access grants, audit log.
  \item \textbf{Container runtime:} Docker + gVisor (\texttt{runsc}) for syscall
    isolation and resource limits.
  \item \textbf{Reverse proxy:} Caddy---automatic TLS, dynamic per-app subdomain
    routing.
  \item \textbf{Object storage:} MinIO (S3-compatible, self-hosted) for application
    bundles and build artefacts.
  \item \textbf{Authentication:} Google OAuth.
\end{itemize}

\begin{figure}[h]
\centering
\begin{verbatim}
+--------------------------------------------------+
|           Silo Web App  (Next.js)                |
|  Workspace . Share . Review . Admin Dashboard    |
+----------------------+---------------------------+
                       | REST / JSON
           +-----------+---------------------------+
           |      FastAPI Control Plane             |
           |  Auth . Apps . Versions . ACLs        |
           |  Capability Diff . Audit Log          |
           +-----------+-----------+---------------+
                       |           |
               +-------+----+ +----+--------+
               | PostgreSQL | | Object Store|
               | system of  | | app bundles |
               | record     | | (MinIO)     |
               +-------+----+ +-------------+
                       |
               +-------+--------+
               | Build & Analysis|
               | Pipeline        |
               | infer -> build  |
               | analyse -> diff |
               +-------+--------+
                       |
               +-------+--------+
               | App Runtime    |
               | Docker+gVisor  |
               | Caddy routing  |
               +----------------+
\end{verbatim}
\caption{Silo system architecture}
\end{figure}

\subsection{CI/CD and fast delivery enabler}

GitHub Actions will:
\begin{itemize}[leftmargin=0.7cm]
  \item Automatically build and run tests on every change (unit tests, Semgrep rule
    tests, golden-file tests for the normaliser).
  \item Produce repeatable deployment artefacts to staging.
  \item Enable safe and frequent releases of incremental features throughout the
    15-week project.
\end{itemize}

% -------------------------------------------------------
\section{Evaluation Criterion %
\BvrHint{\ldots measurable and attributable}}

We define success using measurable metrics tied to the core workflow.

\subsection{Primary evaluation metric}

\textbf{Build success rate on the evaluation set:} percentage of the 10 hand-built
evaluation applications that deploy with zero user configuration on the first attempt.

\begin{itemize}[leftmargin=0.7cm]
  \item \textbf{Target:} $\geq$~8 of 10 applications deploy first-attempt;
    all 10 deploy after at most one auto-repair pass.
  \item \textbf{Attribution:} measured from build logs; each failure is categorised
    (port binding, import resolution, missing entry point, etc.) and a fix is
    documented.
  \item \textbf{Evaluation set:} 10 applications built using AI coding tools
    (Lovable, Claude, Cursor), covering diverse use cases---dashboard, data-upload
    tool, calculator, CRUD form, CSV analyser, reporting tool, internal tracker,
    API client, multi-step workflow, and simple game---so the set exercises the full
    range of build inference cases.
\end{itemize}

\subsection{Secondary evaluation metrics}
\begin{itemize}[leftmargin=0.7cm]
  \item \textbf{Capability-diff accuracy:} for each of the 10 evaluation apps, a
    second version with a known added capability (e.g.\ an outbound network call) is
    produced; Silo's classifier must identify it as capability-expanding.
    Target: 10 of 10 correctly classified.
  \item \textbf{Review comprehension:} 5 non-expert participants are shown two
    capability labels---one neutral diff, one expanding diff---and asked which version
    requests new access. Target: $\geq$~4 of 5 correct.
  \item \textbf{Time-to-publish (TTP):} median elapsed time from application upload
    to live, shareable URL. Target: $\leq$~2 minutes.
  \item \textbf{Cold-start latency:} p50 and p95 from idle container to first HTTP
    response.
  \item \textbf{Collaboration completion rate:} percentage of participants who
    successfully complete the full loop---open a shared app, remix it, make a change,
    submit it, read the capability diff, and propose it upstream.
\end{itemize}

\subsection{Pilot validation plan %
\BvrHint{\ldots high-level}}
\begin{itemize}[leftmargin=0.7cm]
  \item 10 evaluation applications are built using AI coding tools in week~12,
    covering the use cases listed above. Applications are fixed at that point---no
    changes after evaluation begins.
  \item Capability diff test pairs are produced independently: for each app, a second
    version with one known added capability is created and labelled before Silo's
    classifier runs.
  \item Recruit 5 non-expert participants (from non-CS departments) for the review
    comprehension study in week~14.
  \item Metrics are recorded from system logs and study responses; no post-hoc
    adjustments.
\end{itemize}

% -------------------------------------------------------
\section{Scalability %
\BvrHint{\ldots with foresight}}

The proposition should scale in both theoretical and practical senses.

\begin{enumerate}[leftmargin=0.7cm]
  \item \textbf{Scalable (theoretically):} The system supports growth in deployed
    applications and concurrent users by:
    \begin{itemize}[leftmargin=0.7cm]
      \item decoupling build/analysis workers from the API server,
      \item using immutable versioned artefacts content-addressed by \texttt{sha256},
      \item stateless API design with PostgreSQL as the single source of truth.
    \end{itemize}

  \item \textbf{Operationally deployable (practically):} The system runs on a single
    VM:
    \begin{itemize}[leftmargin=0.7cm]
      \item PostgreSQL for persistent state,
      \item Docker + gVisor for container isolation,
      \item Caddy for automatic TLS and subdomain routing,
      \item GitHub Actions for CI/CD to staging and production.
    \end{itemize}
    Scale-to-zero for idle containers prevents resource exhaustion as application
    count grows.
\end{enumerate}

% -------------------------------------------------------
\section{Engine availability heuristic}

A mature set of ``engines'' (tooling/services) is available to implement this as a
web-based project:

\begin{itemize}[leftmargin=0.7cm]
  \item \textbf{Container runtime and isolation:} Docker + gVisor (\texttt{runsc}).
  \item \textbf{Reverse proxy / TLS:} Caddy (automatic certificate provisioning,
    dynamic subdomain routing).
  \item \textbf{Static analyser:} Semgrep (rule-based, deterministic, actively
    maintained).
  \item \textbf{AST parser:} tree-sitter (language-agnostic; shared parse pass for
    build inference and capability analysis).
  \item \textbf{Backend framework:} FastAPI + Pydantic.
  \item \textbf{Frontend framework:} Next.js + TypeScript + Tailwind.
  \item \textbf{Database:} PostgreSQL.
  \item \textbf{Object storage:} MinIO (S3-compatible, self-hosted).
  \item \textbf{Authentication:} Google OAuth.
  \item \textbf{CI runner and staging:} GitHub Actions.
\end{itemize}

We will use these mature components to focus on workflow design, collaboration model
correctness, and measurement---rather than inventing new infrastructure.

% -------------------------------------------------------
\section{Project scope and deliverables %
\BvrHint{\ldots iteration-friendly}}

\subsection{Initial deliverable %
\BvrHint{\ldots first iteration (Weeks 1--7)}}

\begin{itemize}[leftmargin=0.7cm]
  \item Build inference for Python applications: language detection, import resolution,
    entry-point detection, port auto-repair, bounded repair loop.
  \item End-to-end deployment: code upload $\to$ container $\to$ live subdomain URL.
  \item Application workspace: name, description, owner, current version, capability
    summary.
  \item Capability analysis pipeline: Semgrep rules $\to$ normalised manifest $\to$
    nutrition label.
  \item Basic sharing: viewer\,/\,user\,/\,forker\,/\,editor roles, link-based
    access.
  \item User authentication: Google OAuth.
  \item CI: build pipeline + golden-file tests for normaliser + Semgrep rule tests.
  \item CD to staging with a single pipeline job.
\end{itemize}

\begin{table}[h]
\centering
\begin{tabular}{@{}p{1.8cm}p{5.8cm}p{5.6cm}@{}}
\toprule
\textbf{Weeks} & \textbf{Milestone} & \textbf{Exit criterion} \\
\midrule
1
  & Project setup, auth, DB schema, CI/CD
  & Login works; staging deploys on push \\
2--3
  & Build inference (Python): detect, install, run
  & 5 test apps deploy without config \\
4--5
  & Capability analysis + nutrition label
  & Manifest produced; label shown to owner \\
6--7
  & Share dialog + recipient view
  & End-to-end loop demo on staging \\
\bottomrule
\end{tabular}
\caption{Initial deliverable schedule}
\end{table}

\subsection{Subsequent deliverables (Weeks 8--15)}

\begin{itemize}[leftmargin=0.7cm]
  \item Immutable versions + mutable alias routing: stable URL always points to
    approved version.
  \item Capability diff + classification + owner approval flow before any version
    goes live.
  \item Remix/fork workflow: independent copy with its own URL and permissions.
  \item Propose-upstream workflow: forker submits a change; owner reviews capability
    diff; approves or rejects.
  \item Runtime enforcement: egress allowlist from approved manifest, gVisor syscall
    policy, CPU/memory caps.
  \item Audit log: every access, approval, and revoke---timestamped in PostgreSQL.
  \item Org/admin dashboard: live applications, permissions, pending reviews,
    one-click revoke.
  \item Rollback: restore any previous approved version without rebuilding.
  \item Scale-to-zero for idle containers.
  \item Evaluation against 10-app set; review comprehension user study.
  \item Node.js application support (stretch, week~15).
\end{itemize}

\begin{table}[h]
\centering
\begin{tabular}{@{}p{1.8cm}p{5.8cm}p{5.6cm}@{}}
\toprule
\textbf{Weeks} & \textbf{Milestone} & \textbf{Exit criterion} \\
\midrule
8--9
  & Immutable versions + alias routing + capability diff + approval
  & Owner approves diff; URL promotes to new version \\
10--11
  & Remix + propose-upstream workflow + rollback
  & Full fork-and-propose loop demonstrated end-to-end \\
12--13
  & Runtime enforcement + org dashboard + audit log
  & Blocked egress confirmed; admin view shows all apps \\
14
  & Scale-to-zero + evaluation (10-app set + user study)
  & All metrics recorded; study complete \\
15
  & Report, documentation, demo video, buffer
  & Final submission \\
\bottomrule
\end{tabular}
\caption{Subsequent deliverables schedule}
\end{table}

% -------------------------------------------------------
\section{Risks and mitigations}

\begin{itemize}[leftmargin=0.7cm]
  \item \textbf{Build inference fails for unusual frameworks:} Limit initial support
    to Python/Flask and Python/FastAPI; measure build success explicitly; produce
    structured diagnostics on failure rather than an opaque error.
  \item \textbf{Capability normalisation granularity:} if \texttt{/tmp/a.csv} and
    \texttt{/tmp/b.csv} do not collapse to the same canonical capability value, every
    push classifies as privilege-expanding and the diff feature becomes useless.
    Mitigation: golden-file tests on the normaliser from week~6, before evaluation
    apps are built.
  \item \textbf{Evaluation set too easy:} if all 10 apps use identical patterns, the
    evaluation does not exercise edge cases. Mitigation: deliberately span diverse use
    cases (network-heavy, file-heavy, subprocess, no external access) when building
    the set.
  \item \textbf{Single-VM resource exhaustion:} each deployed app consumes memory;
    the VM becomes unusable without scale-to-zero. Mitigation: scale-to-zero is
    scheduled for week~14, before the evaluation window.
  \item \textbf{Collaboration scope creep:} advanced social features can expand
    indefinitely. Mitigation: prioritise the essential loop---remix $\to$ edit
    $\to$ propose $\to$ review $\to$ promote---and defer everything else.
  \item \textbf{Review comprehension study recruitment:} 5 non-expert participants
    needed. Mitigation: recruit from non-CS departments in week~14; study is short
    and structured.
\end{itemize}

% -------------------------------------------------------
\section{Summary}

AI has made the creation of small applications dramatically easier, but the surrounding
software lifecycle remains unnecessarily technical. Silo addresses this gap by providing
a platform where applications can be published without manual deployment configuration,
shared through simple permissions, used directly from a persistent URL, remixed by other
users, and updated through a capability-aware approval workflow that lets owners make
informed decisions without reading source code.

The project meets the course criteria by emphasising time-to-value through a working
end-to-end deployment loop in the first iteration, implementing CI/CD with falsifiable
evaluation metrics (build success rate and capability-diff accuracy against a hand-built
10-app evaluation set), and remaining focused on engineering workflow and system
correctness rather than research-driven novelty. The core claim---that a deterministic
static analyser can produce a human-readable capability label, and that label can gate
version promotion---is testable, measurable, and deliverable within 15~weeks.

\bigskip
\noindent\textbf{The goal is simple: if AI can make software as easily as people make
documents, there should be a way to share and work on that software with the same ease.}

\end{document}
